% ============================================================================
%  THE COLLAPSE MECHANISM -- drop-in section for the ASIA CCS / ACM CCS
%  paper. Split out of what was originally the front half of
%  VLI_method_en.tex, per instruction: Method should describe \ourmethod{}
%  only, not the baselines. This section is contribution (1) of the paper
%  (the proof that aux-using methods collapse); VLI_method_en.tex is
%  contribution (2)/(3) (\ourmethod{}'s two modes).
%
%  Ported and trimmed from TWO sources, deliberately NOT uniformly from
%  VLI_report.tex:
%    - RLU's diagnosis, Prop gap, Prop jensen, Prop lip: from
%      VLI_report.tex ("The failure mechanism"), which is correct for
%      these.
%    - ZLG / ZLG+ / LLG+'s diagnosis: from VLI_baocao_final.tex's
%      "Lo duy nhat con lai" / "Luong nan" argument instead, because
%      VLI_report.tex's Theorem collapse applies Lemma~\ref{lem:norm}
%      (a bias-gradient bound) directly to these three methods, which in
%      the real implementation read the WEIGHT-gradient branch
%      $\widetilde G$ (Section~\ref{sec:prelim-notation}), not $G$. That
%      conflation was caught and fixed in VLI_preliminaries_en.tex; this
%      file carries the fix through to the theorem that depends on it.
%
%  Proof-status convention carried over from the Vietnamese reports:
%  [proved] = short algebraic proof, in Appendix~\ref{app:proofs};
%  [argued] = reasoning given, explicitly NOT a proof -- flagged inline,
%  never upgraded silently; [open] = not yet measured/derived, flagged as
%  a concrete question rather than asserted either way. Do not remove
%  these tags when editing; they are load-bearing for how strong a claim
%  this section is allowed to make.
%
%  Dependencies: same theorem/macro environments as
%  VLI_preliminaries_en.tex, plus \newtheorem{remark}{Remark}.
%  \ref{sec:prelim}, \ref{sec:prelim-notation}, \ref{sec:prelim-facts},
%  \ref{sec:prelim-baselines}, \ref{sec:related}, \ref{sec:eval},
%  \ref{sec:method} (VLI_method_en.tex) point at sections defined
%  elsewhere; \label{app:proofs} is bound in the appendix.
%
%  NOT ported here: numerical verification tables (Section numerics of
%  VLI_report.tex) are evaluation, not mechanism; they belong in the
%  paper's Evaluation section.
% ============================================================================

\section{The Collapse Mechanism}
\label{sec:collapse}

Section~\ref{sec:prelim} gave two facts (Lemma~\ref{lem:norm},
Lemma~\ref{lem:identities}) and five baseline estimators in unified
notation. This section uses the two facts to explain, for each baseline,
\emph{why} it degrades as the attacked model trains -- and in one case,
why it does not. Section~\ref{sec:method} then presents \ourmethod{},
which this analysis motivates but does not itself require.

\subsection{A general collapse argument}
\label{sec:method-collapse}

\begin{theorem}[Collapse of externally-calibrated estimators]
\label{thm:collapse}
Let an estimator of normalized class proportions take the form
$\widehat\pi = \Phi(c - \sigma/n)$, where $\sigma$ is the signal the
estimator reads (either $G$ or $\widetilde G$, Section~\ref{sec:prelim-notation}),
$c$ depends only on information measured outside the attacked window, and
$\Phi$ is $L_\Phi$-Lipschitz in a stated norm. Then
\[
 \bigl\|\widehat\pi - \Phi(c)\bigr\| \;\le\; \frac{L_\Phi}{n}\|\sigma\|_1 .
\]
Hence, whenever $\|\sigma\|_1 \to 0$, the estimate converges to $\Phi(c)$:
a vector determined entirely by $c$ and not by the attacked labels.
\end{theorem}

\begin{proof}[\textnormal{[proved]}]
Lipschitz continuity directly gives the bound; see
Appendix~\ref{app:proofs}.
\end{proof}

By Lemma~\ref{lem:norm}, $\|G\|_1 \to 0$ as $\bar\epsilon\to0$
unconditionally -- this is the case for RLU, which reads $G$ directly
(Table~\ref{tab:baselines}). For ZLG, ZLG$^+$, and LLG$^+$, which read
$\widetilde G$ instead, Theorem~\ref{thm:collapse} applies only once
$\|\widetilde G\|_1 \to 0$ is established, which is not immediate: Section
\ref{sec:method-branch} supplies it, under a condition this section makes
explicit rather than assuming away.

\subsection{Why the bias branch shrinks unconditionally}

Lemma~\ref{lem:norm} already gives $\|G\|_1 \le 2n\bar\epsilon$ with no
further hypothesis, because $\nabla_b = \tfrac1B\sum_i r_i$ (writing
$r_i = p_i - e_{y_i}$) does not involve the embedding at all. RLU's $c$ is
$\mathrm{shift}\cdot\bar\rho$ (Section~\ref{sec:prelim-baselines}), so
Theorem~\ref{thm:collapse} applies to it immediately: $\widehat N \to
n\,\mathrm{shift}/\|\mathrm{shift}\|_1$, a vector fixed by the auxiliary
estimate $\mu_{\mathrm{aux}}$ alone.

\subsection{The weight branch: a dilemma, not a clean bound}
\label{sec:method-branch}

The final layer's weight gradient is $\nabla_{W_{jd}} = \tfrac1B\sum_i
r_{i,j}\,o_{i,d}$, so summing over $d$ gives $\widetilde G_j =
\tfrac1B\sum_i r_{i,j}\,s_i$ with $s_i = \one^\top o_i \ge 0$
(Section~\ref{sec:prelim-notation}). Writing $w_i = s_i/\bar O$ for the
auxiliary estimate $\bar O$ that ZLG$^+$ and ZLG both divide by:

\begin{lemma}[Weight-branch norm bound]
\label{lem:tildenorm}
\[
 \frac{\widetilde G_j}{\bar O} \;=\;
 \underbrace{\frac1n\sum_i r_{i,j}}_{\text{what the estimator assumes}}
 \;+\;
 \underbrace{\frac1n\sum_i (w_i-1)\,r_{i,j}}_{\text{residual}},
 \qquad
 \Bigl\|\frac{\widetilde G}{\bar O}\Bigr\|_1 \;\le\; 2\bar\epsilon(1+V),
\]
where $V = \bigl(\sum_i |w_i-1|\,\epsilon_i\bigr)/\bigl(\sum_i\epsilon_i\bigr)$
measures the error-weighted deviation of $w_i$ from $1$.
\end{lemma}

\begin{proof}[\textnormal{[proved]}]
Split the sum as shown and bound each term; see Appendix~\ref{app:proofs}.
\end{proof}

\begin{remark}[The dilemma]
\label{rem:dilemma}
Lemma~\ref{lem:tildenorm} does not by itself say whether ZLG$^+$ survives,
because it exposes a genuine fork, not a loophole:
\begin{itemize}
 \item If $V$ is small, the bound is tight and $\|\widetilde G/\bar O\|_1
 \to 0$ exactly as $\|G\|_1$ does, so Theorem~\ref{thm:collapse} applies
 to ZLG$^+$ by the same mechanism as RLU. But $V$ small is precisely
 ZLG$^+$'s \emph{own} modeling assumption -- the identity
 $\widehat N_j = n(\bar p_j - \widetilde G_j/\bar O)$
 (Section~\ref{sec:prelim-baselines}) is itself only accurate when
 $s_i \approx \bar O$, i.e.\ when $V$ is small. The estimator's
 justification and its failure mode are the same condition.
 \item If $V$ is large, the residual term does not track the true labels
 -- it is an $(w_i-1)$-weighted sum of a quantity ($r_{i,j}$) that is
 already shrinking, not a label-correlated signal in its own right -- so
 it adds noise to a vanishing estimate rather than rescuing it.
\end{itemize}
The second branch is \textbf{[argued]}, not proved: we have not shown the
residual is literally uninformative, only that nothing in its
construction ties it to the window's labels. Closing this gap -- bounding
how much instance-recovery the residual term can contribute when $V$ is
large -- is open.
\end{remark}

\begin{remark}[LLG$^+$'s calibration ratio may run the other way]
\label{rem:llgp-protect}
LLG$^+$'s calibration ratio is
$\rho \propto (\bar s_{\mathrm{aux}}/\bar s_{\mathrm{train}})\,
(\bar\epsilon_{\mathrm{aux}}/\bar\epsilon_{\mathrm{train}})$: the two
factors move in \emph{opposite} directions as the model overfits, so the
embedding-norm gap that this section identifies as a liability for
ZLG$^+$ could instead partially offset Proposition~\ref{prop:gap}'s
generalization gap for LLG$^+$. This is \textbf{[open]} -- a concrete,
measurable quantity, not yet measured -- and would explain why the two
methods need not degrade identically even though both read $\widetilde
G$.
\end{remark}

ZLG (the original, dummy-data form) is not covered by either remark as
written, since $V$ there is defined relative to a real auxiliary error
rate that a fixed dummy image does not have. Whether ZLG's own $\bar
p,\bar O$ drift away from the window's as training progresses is a
separate, \textbf{[open]} question -- Theorem~\ref{thm:collapse} still
applies in form (its hypothesis is only that $c$ is measured outside the
window, Section~\ref{sec:prelim-baselines}), but the mechanism that would
make $c$ fail to track the window is not Proposition~\ref{prop:gap}'s
generalization-gap argument, which presumes a real data distribution.

\subsection{Why the external magnitude does not shrink with the signal}

\begin{proposition}[Held-out error gap]
\label{prop:gap}
If an example is misclassified by the $\argmax$ rule, then $p_{i,y_i} \le
1/2$ and hence $\epsilon_i \ge 1/2$. Consequently, if the auxiliary error
rate is $\tau_{\mathrm{err}}$, $\bar\epsilon_{\mathrm{aux}} \ge
\tau_{\mathrm{err}}/2$. If $\bar\epsilon_{\mathrm{train}} \to 0$ while
$\tau_{\mathrm{err}}$ stays bounded away from zero, then
$\bar\epsilon_{\mathrm{aux}}/\bar\epsilon_{\mathrm{train}} \to \infty$.
\end{proposition}

\begin{proof}[\textnormal{[proved]}]
See Appendix~\ref{app:proofs}.
\end{proof}

This is the second half of the mechanism, and it applies to RLU, ZLG$^+$,
and LLG$^+$ alike (all three read a real auxiliary error rate): Lemma
\ref{lem:norm}/\ref{lem:tildenorm} say the attacked signal shrinks;
Proposition~\ref{prop:gap} says the auxiliary magnitude entering $c$ does
not. Their ratio diverges, so comparing them is a comparison between two
quantities on incommensurable scales -- \emph{generalization gap, not
estimator design, is what breaks these attacks.}

\subsection{Two further defects specific to RLU}

\begin{proposition}[Jensen gap in logit averaging]
\label{prop:jensen}
For a sample with margin $M$, $\epsilon = (1+e^M)^{-1} \approx e^{-M}$ for
$M$ large and positive. If within-class margins are approximately
$\mathcal N(\mu,\sigma^2)$, $\E[e^{-M}] = e^{-\mu+\sigma^2/2} \ge
e^{-\E[M]} = e^{-\mu}$ by Jensen's inequality, with equality only at
$\sigma=0$.
\end{proposition}

RLU's $\mathrm{shift} = \softmax(\mu_{\mathrm{aux}})$ substitutes a
softmax of \emph{mean logits} for the mean of softmaxes -- exactly the
substitution Proposition~\ref{prop:jensen} describes, systematically
\emph{understating} error mass by a factor growing with $\sigma^2$, which
training increases: the bias is largest precisely when the signal is
weakest.

\begin{proposition}[Clip-and-renormalize is not Lipschitz; projection is]
\label{prop:lip}
Let $R(y) = n\pos{y}/\|\pos{y}\|_1$ (used by RLU and, after clipping, by
ZLG$^+$) and $\Pi_\simplex$ denote Euclidean projection onto $\simplex$.
Then $R$ has unbounded Lipschitz constant while $\Pi_\simplex$ is
non-expansive everywhere.
\end{proposition}

\begin{proof}[\textnormal{[proved]}]
See Appendix~\ref{app:proofs}.
\end{proof}

RLU's final normalization can therefore \emph{amplify} whatever noise
survives the collapse above, rather than merely fail to correct it;
Section~\ref{sec:eval} reports a measured amplification factor for $R$
against exactly $1$ for $\Pi_\simplex$.

\subsection{iLRG: a conditioning failure, not a calibration failure}

iLRG uses no auxiliary statistic, so Theorem~\ref{thm:collapse} does not
apply to it. Its difficulty is different: by Lemma~\ref{lem:identities},
$G_j = -a_j + E_j$, so both terms vanish as $\bar\epsilon\to0$, and the
bias gradient used as a \emph{denominator} in $\hat e_j =
\nabla_{W_j}/\nabla_{b_j}$ tends to zero. The ratio's relative error is
amplified without bound even though it remains finite in exact
arithmetic, and the subsequent linear system grows ill-conditioned as
$p_i\to e_{y_i}$ collapses every recovered probability vector toward a
one-hot vector.

\subsection{Summary of baseline diagnoses}

\begin{table*}[t]
\centering
\caption{How each baseline fails (or does not), and the status of each
claim.}
\label{tab:diagnoses}
\begin{tabular}{@{}lp{3.6cm}p{6.8cm}l@{}}
\toprule
Method & Mechanism & Limit or symptom & Status \\
\midrule
iLRG & conditioning & division by a vanishing bias gradient; near-singular pseudo-inverse & proved \\
LLG & none (scale-free) & survives; assumes $r_j\equiv1$, uniform $Q$, multiplies a single-batch answer by $m$ & proved \\
ZLG & Thm~\ref{thm:collapse}; branch condition in Remark~\ref{rem:dilemma} & $\widehat N \to n\bar p$ if $V$ small; mechanism for $c$ drifting is open for dummy data & mixed \\
LLG$^+$ & Thm~\ref{thm:collapse} + Prop~\ref{prop:gap}; possibly offset by Remark~\ref{rem:llgp-protect} & auxiliary offset dominates $\widetilde G$; net effect on LLG$^+$ vs.\ LLG is open & mixed \\
ZLG$^+$ & Thm~\ref{thm:collapse} + Prop~\ref{prop:gap}; branch condition in Remark~\ref{rem:dilemma} & $\widehat N \to n\bar p \approx$ uniform prior & mixed \\
RLU & Thm~\ref{thm:collapse} + Props~\ref{prop:gap},~\ref{prop:jensen},~\ref{prop:lip} & $\widehat N \to n\,\mathrm{shift}$; Jensen bias $e^{\sigma^2/2}$; non-Lipschitz normalization & proved \\
\bottomrule
\end{tabular}
\end{table*}
